%%%%%%%%%% Section 5 %%%%%%%%%%
\vspace{-0.3ex}
\section{Policy Construction and Execution}
\label{sec-classes}
\vspace{-0.2ex}

This section develops goal-aware policies for the
inference agent and an algorithm for their adaptive
execution. Section~\ref{sec-policies} constructs
score-based, learned, and hybrid policies using the
scores of Section~\ref{sec-score}.
Section~\ref{sec-execution} presents the %agent's
deduction
algorithm and shows its soundness,
goal-level completeness, and termination guarantees.
%\looseness=-1

%%%%%%%%%% Section 5.1 %%%%%%%%%%
\vspace{-1ex}
\subsection{Policy Construction}
\label{sec-policies}
\vspace{-0.2ex}

A scoring function ranks individual inference
actions, while a policy determines which action
the agent evaluates next. We first construct
policies directly from $f(S,g,\mu)$, and then
develop learned policies that exploit previous
executions to estimate downstream costs.
All policies operate on the available actions
$\mathcal A(S)$ of Section~\ref{sec-model}.

\stitle{Score-based policies.}
A {\em deterministic policy} selects the  action that has the
highest score, \ie we select action $\mu$ such that
\vspace{-0.7ex}\[
\pi_{\rm det}(S,g)=
\arg\max_{\mu\in\mathcal A_g(S)}f(S,g,\mu),
\vspace{-0.7ex}\]
where
$\mathcal A_g(S)=
\{\mu\in\mathcal A(S)\mid r(g,\mu)>0\}$
contains the goal-relevant available actions.
Ties are resolved by a fixed order.
The policy favors actions with stronger
symbolic relevance and supporting evidence,
but lower estimated evaluation costs.

Since $f$ is an estimate, the score may favor an action that turns out
to be costly or unsuccessful; thus the agent should
also explore other promising alternatives.
A {\em randomized policy} does so by selecting actions
with probabilities proportional to their scores:
\vspace{-0.7ex}\[
\pi_{\rm rnd}(\mu\mid S,g)=
\frac{\exp(f(S,g,\mu)/\tau)}
{\sum_{\mu'\in\mathcal A_g(S)}
\exp(f(S,g,\mu')/\tau)},
\vspace{-0.7ex}\]
where $\exp(x)=e^x$ is the exponential function,
assigning higher selection probabilities
to higher-scoring actions, and $\tau>0$ controls
this preference. Small $\tau$ favors the
highest-scoring action, and larger $\tau$
promotes exploration. Randomization reduces
reliance on potentially inaccurate scores,
but does not alone guarantee completeness.
\looseness=-1

\begin{example}
\label{exa-policy}
Consider the AML goal
$g=\textsf{HighRisk}(v)$ of
Example~\ref{exa-chase}.
After action $\mu_1=(\varphi_1,h_1)$ derives
$\textsf{HighRisk}(u)$, both
$\mu_2=(\varphi_2,h_2)$ and
$\mu_3=(\varphi_3,h_3)$ can derive the goal $g$.
If previously validated transfer conditions
give $\mu_2$ a higher score, then the deterministic
policy selects $\mu_2$. In contrast,
the randomized policy favors $\mu_2$ but
may explore $\mu_3$, which uses shared
infrastructure instead.
\end{example}
\vspace{-1.4ex}

\stitle{Learned policies.}
The handcrafted score considers immediate
relevance, evidence, and evaluation cost,
but cannot fully capture the downstream
computational consequences of selecting an
action. We thus learn from execution
experience.
Let $\Phi(S,g,\mu)$ denote the features of
a candidate action, including its symbolic
relevance $r(g,\mu)$, evidence estimate
$\widehat e(S,\mu)$, remaining evaluation-cost
estimate $\widehat c_{\rm eval}(S,\mu)$,
and the representations $\mathbf v_\mu$
and $\mathbf E(S)$ from Section~\ref{sec-score}.

Rather than predicting whether an action
appears in a successful derivation, we learn
its {\em action-conditioned cost-to-go}.
Let $\pi_b$ be a continuation policy used
to collect training executions.
Define
\vspace{-0.7ex}\[
Q^{\pi_b}(S,g,\mu)=
\mathbb E\left[
\cost(C)\mid S,\mu,\pi_b
\right],
\vspace{-0.7ex}\]
where $C$ starts by evaluating $\mu$ at state $S$
and then follows policy $\pi_b$ until reaching $g$.
Here $\mathbb E$ averages the execution costs over
the possible action sequences induced by $\pi_b$,
since randomized choices may lead to different
costs. For a deterministic $\pi_b$, the expectation
reduces to the cost of its single execution.
The cost includes all action evaluations and
policy overhead, as in PSP. An execution that
fails to reach a derivable goal has infinite cost.


To implement this policy,
we train a model $\widehat Q_\theta$
to estimate this cost from
$\Phi(S,g,\mu)$.
For training, we collect completed executions
under $\pi_b$, retaining the state, selected
action, and remaining observed execution cost
at each decision.
Additional rollouts from recorded states
evaluate alternative actions and their
subsequent executions. These observations
provide action-conditioned cost targets;
unobserved alternatives are not assigned
artificial costs.


For each training observation $(S,g,\mu,y)$,
let $y$ be the observed remaining execution cost
after selecting $\mu$ at state $S$ and following
$\pi_b$. We train a model
$\widehat Q_\theta(S,g,\mu)$, parameterized by
$\theta$, to predict this cost from the features
$\Phi(S,g,\mu)$. The model minimizes
\vspace{-0.7ex}\[
\mathcal L(\theta)=
\sum_{(S,g,\mu,y)\in\mathcal D}
\big(\widehat Q_\theta(S,g,\mu)-y\big)^2,
\vspace{-0.7ex}\]
where $\mathcal D$ contains the collected
finite-cost training observations. Note that
failed goal-reaching rollouts are recorded
separately rather than assigned finite targets.
During execution, the agent combines learned
selection with {\em fair} exploration: it periodically
selects the longest-waiting available action,
ensuring that no candidate is indefinitely
ignored because of a low predicted score.
Together with the \revise{revival} of dynamically
suspended actions, this preserves goal-level
completeness (see Theorem~\ref{thm-complete} shortly).



Unlike immediate scoring, cost-to-go learning
can recognize that a seemingly inexpensive
action may trigger costly subsequent deductions.
However, $\widehat Q_\theta$ estimates costs
under the continuation policy $\pi_b$;
it does not guarantee globally optimal
policy selection.


\stitle{Hybrid neuro-symbolic policies.}
Learned policies may assign high scores to
actions that are unlikely to contribute to
the target goal. To prevent such choices,
we combine symbolic reasoning with learned
action selection in two stages.

First, we use the rule dependency graph
to restrict the available actions to
$\mathcal A_g(S)$.
By Proposition~\ref{prop-relevance},
excluding actions with $r(g,\mu)=0$
preserves goal derivability.
Second, among the remaining candidates,
the agent selects the action with the
lowest predicted cost-to-go
$\widehat Q_\theta(S,g,\mu)$.
Alternatively, it can rank candidates
using the handcrafted score $f(S,g,\mu)$.
Thus, symbolic filtering identifies
potentially useful deductions, while
learned guidance estimates which action
can reach the goal at lower computational
cost.

We also support LLM-guided selection,
\revise{instantiated with the decision model
Jev~\cite{Jev}.}
Given the target goal $g$, relevant
\GARs, grounded candidate actions,
and the current execution evidence,
\revise{Jev} ranks actions in $\mathcal A_g(S)$
according to their expected contribution
to the goal and evaluation cost\revise{,
returning a calibrated ranking over the
candidates in one pass~\cite{Jev-benchmark}}.
The agent selects a recommended action
and validates its preconditions using
the exact \GAR semantics.
The resulting facts and evaluation outcomes
update the deduction state for subsequent
decisions.

In both cases, learned models and LLMs
control only the order of inference actions, without
introducing unsupported facts.
Combined with fair selection and
retrial of suspended actions,
this preserves soundness and goal-level
completeness.



%%%%%%%%%% Section 5.2 %%%%%%%%%%
\vspace{-1ex}
\subsection{Policy-Guided Deduction}
\label{sec-execution}
\vspace{-0.2ex}

\revise{We integrate the policies into an adaptive
deduction algorithm.
\AdaDeduce derives the goal in two loops:
a {\em main loop} that drives the deduction
toward the goal under the policy, and a
{\em fixpoint loop} that completes the
inference steps that had to wait for
premises still to be derived.
It stops when the goal is derived, or when
no further fact can be derived.}

\vspace{-0.4ex}
\stitle{Working sets.}
\revise{Recall from Section~\ref{sec-model}
the deduction state $S=(G,\Sigma,\Gamma,H)$.
The agent maintains two collections on top
of it.
The {\em ready set} $\mathcal R$ consists of
pairs $(\varphi,h)$, where $h$ is a complete
match of the pattern of a \GAR
$\varphi$; these are the discovered
goal-relevant actions awaiting evaluation.
The {\em suspended set} $\mathcal W$
consists of such pairs annotated with the
dynamic preconditions still missing,
together with the preconditions already
validated; these are the partial results
that wait for their missing facts.
%$\mathcal W$ is managed by the algorithm
%alone and never competes for selection.
The two sets exchange objects in a closed
loop: graph matching instantiates
goal-relevant rules into actions, which
enter $\mathcal R$; predicate evaluation
then checks the preconditions of an action
of $\mathcal R$ under its match, deriving
its consequent into $\Gamma$,
or suspending it in $\mathcal W$
when only dynamic preconditions are
missing; and new facts complete suspended
actions.}

\vspace{-0.2ex}
\stitle{Agent.}
\revise{We implement \AdaDeduce in
Algorithm~\ref{alg-policy}.
Given a graph $G$, \GARs $\Sigma$, initial
facts $\Gamma_0$, and target goal $g$, the
agent builds $D_{\Sigma,g}$ once (line 1),
marking each node as a {\em leaf} if it has
no incoming edge, and {\em derivable}
otherwise; the two loops below run on
$D_{\Sigma,g}$ and its marks.
It derives $g$ in the main loop and the
fixpoint loop, and returns whether $g$ is
derived, the resulting facts, the complete
execution trace, and the actual execution
cost.}

\begin{figure}[t]
\centering
\begin{minipage}{0.98\columnwidth}
\small
\myhrule\\[0.5ex]
{\bf Input}: Graph $G$, \GARs $\Sigma$,
initial facts $\Gamma_0$, goal $g$, policy $\pi$.\\
{\bf Output}: Trace $C$, facts $\Gamma$, cost $c$.\\[0.5ex]
\bcc \hspace{2ex}~$\Gamma\gets\Gamma_0$;
$H,\mathcal R,\mathcal W\gets\emptyset$;
$c\gets0$; \revise{build $D_{\Sigma,g}$;}\\
\icc \hspace{2ex}~\revise{\ForEach star $\varphi$
of $D_{\Sigma,g}$ \Do}\\
\icc \hspace{2ex}~\quad\revise{match the pattern
of $\varphi$ in $G$, reusing matches of
shared sub-patterns, and add the actions
$(\varphi,h)$ to $\mathcal R$;}\\
\icc \hspace{2ex}~\quad\While
$\mathcal R\neq\emptyset$ \Do\\
\icc \hspace{2ex}~\qquad\revise{select
$\mu\in\mathcal R$ by $\pi$;}\\
\icc \hspace{2ex}~\qquad\revise{check the
preconditions of $\mu$, reusing outcomes
in $H$;}\\
\icc \hspace{2ex}~\qquad\If all preconditions
hold \Then $\Gamma\gets
\Gamma\cup\{h(p_0)\}$, where $\mu=(\varphi,h)$;\\
\icc \hspace{2ex}~\qquad\revise{charge
$c_{\rm sel}$ and $c_{\rm eval}$ to $c$;
\If $h(p_0)=g$ \Then \Return $(C,\Gamma,c)$;}\\
\icc \hspace{2ex}~\qquad\Else \If a leaf
precondition $p$ fails \Then \revise{drop
$\mu$, record $p$ in $H$, and prune all
actions sharing $p$;}\\
\icc \hspace{2ex}~\qquad\Else \revise{suspend
$\mu$ in $\mathcal W$ with its validated
and missing preconditions;}\\
\icc \hspace{2ex}~\qquad\revise{$S'\gets(G,\Sigma,\Gamma,H)$;
$C\gets C\circ\langle(S,\mu,S')\rangle$;}\\
\icc \hspace{2ex}~\revise{\Repeat}\\
\icc \hspace{2ex}~\quad\revise{\ForEach
$\mu\in\mathcal W$ \Do}\\
\icc \hspace{2ex}~\qquad\revise{\If the missing
preconditions of $\mu$ are all in $\Gamma$
\Then apply $\mu$ as in lines 7--8 and 11,
and remove it from $\mathcal W$;}\\
\icc \hspace{2ex}~\qquad\revise{\Else \If some
missing precondition of $\mu$ has no
remaining action deriving it \Then drop
$\mu$;}\\
\icc \hspace{2ex}~\revise{\Until
$\mathcal W=\emptyset$ or a round derives
no new fact;}\\
\icc \hspace{2ex}~\Return $(C,\Gamma,c)$.\\[0.5ex]
\myhrule
\end{minipage}
\vspace{-2.4ex}
\caption{\revise{Agent \AdaDeduce: policy-guided
deduction with star expansion and fixpoint
revival.}}
\label{alg-policy}
\vspace{-4.3ex}
\end{figure}

\etitle{The main loop.}
\revise{The agent works in rounds over
$D_{\Sigma,g}$, one star per round, each
round in two stages: the generation stage
produces the candidate actions of the star,
and the policy stage evaluates them before
the next round begins (lines 2--11).}

\eetitle{(1) The generation stage.}
\revise{This stage produces all candidate
actions of the rules in $D_{\Sigma,g}$.
It expands one {\em star} of $D_{\Sigma,g}$ at a time,
a rule with its consequent as the center
and its preconditions as the spokes
(Figure~\ref{fig-dSigma}), by matching
its pattern in $G$; each complete match $h$
forms one candidate action $(\varphi,h)$
and enters $\mathcal R$ (lines 2--3).
Matching checks vertex labels and edges
only, and shared
sub-patterns are matched once and reused.
%Since $G$ never changes, this stage
%determines the candidate actions of the
%entire execution.
}

\vspace{-0.2ex}
\eetitle{(2) The policy stage.}
\revise{As the downstream stage of each
generation stage, it turns the candidate actions of
$\mathcal R$ into facts under the control
of the policy $\pi$.
Whenever $\mathcal R$ is nonempty,
the agent evaluates its actions until
$\mathcal R$ turns empty (lines 4--11): the
policy picks the next action (line 5), and
the agent checks its preconditions, \ie predicates, 
reusing
outcomes in $H$ (line 6).
The marks of $D_{\Sigma,g}$ split the
actions into two kinds
(Section~\ref{sec-model}): a {\em static}
action has only leaf spokes and is settled
by one evaluation; a {\em dynamic} action
has a derivable spoke and may have to wait.
Three outcomes are possible.
(a) All preconditions hold: the action
derives its consequent, $\Gamma$ is
updated, and the cost is charged, stopping
at once if the fact is $g$ (lines 7--8).
(b) A leaf precondition fails: the action
can never succeed and is discarded, with
every action sharing that failed
precondition (line 9).
(c) Only derivable facts are missing: the
action is suspended in $\mathcal W$ with
its bindings and validated preconditions
(line 10).
In each case the transition is recorded in
the trace (line 11), and the agent proceeds
to the next action of $\mathcal R$.}

\vspace{-0.2ex}
\etitle{The fixpoint loop.}
\revise{When the main loop ends,
$\mathcal W$ may still hold partial results
waiting for missing facts, and the fixpoint loop
closes them out in rounds (lines 12--13).
No matching is needed anymore: all matches
were computed in the main loop, and
$G$ never changes; each round only looks up
the missing facts of $\mathcal W$ in $\Gamma$, reusing the
recorded partial results.
An action whose missing facts have all
arrived succeeds and leaves $\mathcal W$
(line 14), possibly completing others in
later rounds; one waiting on a fact that no
remaining action can derive is abandoned
(line 15).
The rounds iterate until $\mathcal W$ is
empty or $\Gamma$ stays the same for two
consecutive rounds, \ie it reaches a
fixpoint (line 16): only a new fact could
complete a partial result, so nothing in
$\mathcal W$ can ever succeed again.
Together with the early stop at $g$
(line 8), \AdaDeduce terminates, returning
the trace, facts and cost (line 17).}

\begin{example}
\label{exa-execution}
Continue Example~\ref{exa-policy}, with
$g=\textsf{HighRisk}(v)$.
\revise{Suppose that the generation stage
expands the star of $\varphi_2$ first, and
the policy stage picks
$\mu_2=(\varphi_2,h_2)$ from $\mathcal R$.}
Its ownership, transfer, and amount
conditions hold, but
$\textsf{HighRisk}(u)$ is missing.
\revise{Since only a dynamic precondition
fails, the action is suspended}
in $\mathcal W$; its validated
preconditions are recorded with it as a
partial result.
\revise{When the star of $\varphi_1$ is
expanded,} the agent evaluates
$\mu_1=(\varphi_1,h_1)$, deriving
$\textsf{HighRisk}(u)$.
\revise{After all stars are processed, the
fixpoint loop revisits $\mu_2$ in its
first round: the missing precondition is
now in $\Gamma$, the recorded partial
result is reused without re-evaluation,
and $\mu_2$ derives
$g=\textsf{HighRisk}(v)$.}
Execution stops immediately.

A learned policy may
revisit $\mu_2$ before $\mu_3$ in the
fixpoint loop
because its estimated remaining cost is
lower. The earlier unsuccessful evaluation
is thus useful: its evidence reduces
the cost of reaching the goal.
\end{example}
\vspace{-1.4ex}

\stitle{Correctness guarantees.}
We next show that policy-guided deduction
preserves logical validity and derives
every reachable goal under fair execution.
These guarantees hold due to exact \GAR
evaluation and the working-set discipline,
not on scoring accuracy.

We first show the soundness: heuristic, learned, or LLM-based
selection cannot introduce unsupported
deductions.

\begin{theorem}
\label{thm-sound}
For any policy $\pi$, every fact derived
by \AdaDeduce belongs to
$\Gamma^*=\Chase(G,\Sigma,\Gamma_0)$.
\end{theorem}
\vspace{-1ex}

\proofS~We show by induction on the execution steps 
that $\Gamma\subseteq\Gamma^*$.
Initially, $\Gamma=\Gamma_0\subseteq\Gamma^*$.
For the inductive step, note that 
each successful action validates all
preconditions of a \GAR and performs
a sound chase step.
Failed actions add no facts.
\eop

\vspace{0.8ex}
We next show the goal-level completeness, \ie goal reachability even
when the policy's scores are inaccurate.

\begin{theorem}
\label{thm-complete}
For finite $G$ and $\Sigma$, under
complete goal-relevant candidate discovery
and fair execution,
\AdaDeduce derives
%every goal
all $g\in\Gamma^*$.
\end{theorem}\looseness=-1

\vspace{-1ex}

\proofS~We prove this also by induction
on the number of chase derivation steps.
Since $g\in\Gamma^*$, a finite chase
derivation of $g$ exists.
By Proposition~\ref{prop-relevance},
every action needed for this derivation
survives symbolic filtering.
\revise{Every action of the derivation is
generated and evaluated, since the
generation stage expands every star in the
cone of $g$.
Such an action either succeeds or suspends
only on dynamic preconditions; those are
derived by earlier actions of the
derivation, which succeed by induction, and
the fixpoint rounds revisit the action
until they arrive.
The fixpoint cannot end before $g$ is
derived: the first action of the derivation
not yet succeeded would have all its
preconditions in $\Gamma$ and would derive
a new fact.}
\eop

\vspace{0.8ex}
Finally, we show that our algorithm terminates, and
repeated unsuccessful
evaluations cannot cause nontermination.

\begin{theorem}
\label{thm-termination}
For finite $G$ and $\Sigma$,
\revise{\AdaDeduce terminates.}
\end{theorem}%\looseness=-1

\vspace{-1ex}

\proofS~The finite graph admits finitely many
grounded rule applications and ground facts.
\revise{Each star is expanded once, so each
action is generated once, and the policy
stage evaluates each generated action once.
In the fixpoint loop, every round except
the last derives at least one new fact;
since $\Gamma$ grows monotonically and only
finitely many facts exist, the number of
rounds is finite.}
\eop

\etitle{Computational cost.}
\revise{The execution cost of a trace is
accounted as in Section~\ref{sec-statement}.}
By Section~\ref{sec-score}, scoring a
candidate with a precomputed representation
takes $O(D+|P_\mu|)$ time, where $D$ is
the HDC dimension.
\revise{Generation matches only the rules
of $D_{\Sigma,g}$, each star exactly once,
and the fixpoint loop re-checks only the
missing preconditions of suspended actions,
at most $|\mathcal W|$ lookups per round.}
Incremental maintenance of the working sets
avoids repeating unaffected graph matches,
and recorded precondition outcomes avoid
redundant evaluations.
The saving depends on the reduction in
inference work relative to these control
costs, which we measure experimentally.\looseness=-1

\etitle{Remark.}
Unlike Meta-Chase~\cite{Meta-Chase}, which
controls rule-set selection, we
select individual grounded inference
actions to minimize goal-reaching cost.
\revise{Its novelty lies in combining
cost-to-go learning, reusable execution
evidence, and cone-restricted star expansion
over the working sets with sound, complete,
and terminating \GAR deduction.}
